\documentclass[answers,12pt,addpoints]{exam} 
\usepackage{import}

\import{C:/Users/prana/OneDrive/Desktop/MathNotes}{style.tex}

% Header
\newcommand{\name}{Pranav Tikkawar}
\newcommand{\course}{Spatial Statistics}
\newcommand{\assignment}{Homework 4.1}
\author{\name}
\title{\course \ - \assignment}

\begin{document}
\maketitle

\newpage
\begin{questions}
\question What is the trade off for model simplicity and estimation reliability. Is there a metric that can be used to determine the ease of estimation of a model?

\question Also for a model, what is the property call that refers to that with a close number of observations, the model will give similar results? I am pretty sure it is stability, I am not sure if has another name in spatial contexts. Also how necessary or strong does this property need to be for a model to be considered good?

\question Reparameterizing the model to make the likelihood function more interpretable is a really good idea. This is like a change of basis that reorients the model with parameters being as "orthogonal" as possible. I am curious as to what are some of the other methods that can be used to make the likelihood function more interpretable. Especially for time valued data.

\question I dont understand why we need to consider the concept of there existing multiple sequences of MLE $(hat \phi_{1m}, \hat \phi_{2m}, )$ based on the the number of observations $m$. Should the MLE be independent of the sample size but instead it depends on the density of the data in the domain of the model? or are these equivalent? 

\question This chapter honestly was very straight forward. I frankly dont have many questions about it. I do realized that this chapter focused on a very specific type of model, but the questions that I would have about the more general case or even a specific case of a more complex model are not really materialized enough in my mind to ask.

\question My main takeaway from this chapter (broadly speaking) is that reparameterization is a helpful tool not just for interpretability but also for estimation. I also realize that both the fixed-domain and the increasing-domain asymptotic are important to consider when analyzing the properties of a model. 

\end{questions}

\end{document}